
A program that outputs the correct answer on 764 tests can still be wrong 40\% of the time --- if one evaluates it with a formal contract on the inputs that matter most to contract semantics.

This finding motivates the central question of this paper: when the evaluation oracle changes from differential output equality to formal contract satisfaction, how many programs classified as correct by dense testing are exposed as incorrect? The answer is nearly half --- and understanding why requires distinguishing what different oracle types actually measure.

\textbf{The surface problem} is well-known: LLM-generated code frequently contains errors that sparse test suites miss. The field has responded with denser evaluation --- EvalPlus expanded HumanEval by $80\times$ and MBPP by $35\times$, reducing pass@1 by up to 23.1\% [Liu2023EvalPlus]. Yet these exhaustive differential test suites rely on a fundamental design choice: they check whether a program produces the same output as a reference implementation on a finite set of sampled inputs.

\textbf{The deeper problem} is that this oracle cannot check universal properties. Formal contracts --- pre- and post-conditions on function behavior --- encode relational invariants, quantified conditions, and semantic constraints that must hold for all valid inputs, not just a finite sample. A program can return the numerically correct answer on every test input while violating a postcondition asserting, for example, that the output is a permutation of the input or that all returned elements satisfy a membership constraint. Dense differential testing cannot detect these violations by construction.

\textbf{The gap} is this: no published work has measured the execution-based oracle-strength gap across multiple LLM families on ContractEval [Lim2025ContractEval] --- the only benchmark pairing standard coding tasks with formal Python pre/post-condition contracts. ContractEval's own evaluation uses SMT-based test synthesis (Z3), which is tractable for only 25.82\% of tasks. Prior property-based testing work [Bose2025Prompts] evaluates only two models and lacks formal contract annotations. No study has quantified how much stronger a contract oracle is than a dense differential oracle on the same inputs.

\textbf{The key methodological insight} is that oracle semantics must be decoupled from input strategy. By using ContractEval's contract-violating test inputs (CVT inputs) as a controlled, fixed input set, both a differential equality oracle and a contract oracle can be evaluated on identical inputs --- directly measuring oracle semantic strength without any input-generation confound. On CVT inputs, when an LLM program returns the same output as the ground-truth canonical, the differential oracle classifies the program as correct. If the reference contract fires nonetheless, that disagreement is attributable entirely to the contract encoding semantics that equality cannot express.

\textbf{This paper makes the following contributions:}

\begin{enumerate}
\item \textbf{Oracle isolation design} --- a controlled experiment using CVT inputs to measure contract oracle semantic strength independently of input generation strategy; to our knowledge, the first controlled methodology for oracle-type comparison on a benchmark with formal contract annotations.
\end{enumerate}

\begin{enumerate}
\item \textbf{Quantification of the differential-contract oracle gap} --- oracle-isolation gap of 0.40 (95\% CI [0.358, 0.445]) across 364 ContractEval tasks and 3 LLM families (Experiment A; 10,432 triples), with 40\% contract-unique mass; both values exceed predicted thresholds by $4\times$ and $8\times$ respectively.
\end{enumerate}

\begin{enumerate}
\item \textbf{Adaptive PBT contribution} --- icontract-hypothesis adaptive property-based testing adds a mean ~10 percentage points of violations beyond static oracle inputs (Wilcoxon p = 2.64e-22), consistently across all 5 tested model families (Experiment B; 17,226 triples); confirming that adaptive search and oracle semantics contribute distinct, independent detection layers.
\end{enumerate}

\begin{enumerate}
\item \textbf{Principled negative results} --- oracle-isolation gap does not scale monotonically with contract richness (Spearman $\rho$ = 0.136), and cross-model gap is negligible at n = 5 ($\Delta R^2$ = 0.004); both findings are informative about ContractEval's contract structure and the limits of the current evaluation scope.
\end{enumerate}

Section 2 reviews related work. Section 3 presents methodology. Section 4 describes experimental configuration. Section 5 reports results. Section 6 discusses implications and limitations. Section 7 concludes.

